\documentclass[12pt, a4paper]{ctexart}
\usepackage{fontspec}
\usepackage{xcolor}
\usepackage{hyperref}
\usepackage{geometry}
\usepackage{enumitem}
\usepackage{parskip}
\usepackage{titlesec}
\usepackage{tcolorbox}
\tcbuselibrary{breakable}
\usepackage{setspace}
\usepackage{listings}
\usepackage{amssymb}

\geometry{left=2.5cm, right=2.5cm, top=2.5cm, bottom=2.5cm}

\setmainfont{Times New Roman}
\setCJKmainfont{SimSun}

\hypersetup{
    colorlinks=true,
    linkcolor=blue!70!black,
    urlcolor=cyan,
    pdftitle={面向对象上-逐题精讲解义},
    pdfauthor={专升本-fifine老师}
}

\definecolor{myblue}{RGB}{30,100,200}
\definecolor{lightblue}{RGB}{230,240,255}
\definecolor{darkblue}{RGB}{0,50,120}

\newtcolorbox{bluebox}[1][]{
    breakable,
    colback=lightblue,
    colframe=myblue,
    coltitle=darkblue,
    fonttitle=\bfseries,
    title={#1},
    boxrule=0.8pt,
    arc=4pt,
    left=6pt,
    right=6pt,
    top=6pt,
    bottom=6pt,
    before skip=8pt,
    after skip=8pt
}

\usepackage{etoolbox}
\BeforeBeginEnvironment{bluebox}{\par\vfill}%
\AfterEndEnvironment{bluebox}{\par\vfill}%

\titleformat{\section}
  {\normalfont\Large\bfseries\color{myblue}}{\thesection}{1em}{}
\titlespacing*{\section}{0pt}{3.5ex plus 1ex minus .2ex}{1.5ex plus .2ex}

\title{\textcolor{myblue}{\textbf{面向对象（上） - 逐题精讲解义}}}
\author{专升本-fifine老师 教学组}
\date{2026年03月02日}

\lstset{
    language=Java,
    basicstyle=\ttfamily\small,
    keywordstyle=\color{blue},
    commentstyle=\color{green!60!black},
    stringstyle=\color{orange},
    numbers=left,
    numberstyle=\tiny\color{gray},
    frame=single,
    breaklines=true,
    columns=fullflexible,
    showstringspaces=false
}

\newcommand{\code}[1]{\texttt{#1}}

\begin{document}

\maketitle
\thispagestyle{empty}
\newpage

\section{面向对象（上）}

\begin{enumerate}[label=\textbf{\arabic*.}]

	\item \textbf{以下关于类和对象的说法正确的是}
	      \begin{enumerate}[label=\Alph*.]
		      \item 类是对象的实例
		      \item 对象是类的模板
		      \item 一个类只能创建一个对象
		      \item 对象是类的具体实例
	      \end{enumerate}

	      \begin{bluebox}{答案与深度解析}
		      \textbf{答案：D. 对象是类的具体实例}

		      \textbf{【核心认知框架>}
		      \begin{itemize}[label={}, leftmargin=*, nosep]
			      \item \textbf{题目本质}：考查面向对象最基础的概念——类与对象的关系
			      \item \textbf{关键区分点}：类是抽象模板，对象是具体实例；一个类可以创建多个对象
		      \end{itemize}

		      \textbf{【选项原理深度拆解】}
		      \begin{itemize}[label={}, nosep]

		      \item \textbf{类是对象的实例 — 错误（概念颠倒）}
		      \begin{itemize}[label={}, nosep]
			      \item 对象是类的实例（Instance）
			      \item 类不是对象的实例
		      \end{itemize}

		      \item \textbf{对象是类的模板 — 错误（模板是类）}
		      \begin{itemize}[label={}, nosep]
			      \item 类是模板（Template），对象是实例（Instance）
		      \end{itemize}

		      \item \textbf{一个类只能创建一个对象 — 错误（可以创建多个）}
		      \begin{itemize}[label={}, nosep]
			      \item 类是模板，可以反复使用
			      \item 每次new都创建新对象
		      \end{itemize}

		      \item \textbf{对象是类的具体实例 — 正确（正确定义）}
		      \begin{itemize}[label={}, nosep]
			      \item 类：抽象的模板，定义属性和方法
			      \item 对象：类的具体实例，具有实际状态和行为
		      \end{itemize}
		      \end{itemize}

		      \textbf{【应背诵知识点>}
		      \begin{itemize}[label={}, nosep]
			      \item 类是抽象模板，对象是具体实例
			      \item 一个类可以创建多个对象
			      \item 对象拥有类的所有属性和行为
		      \end{itemize}

		      \textbf{【命题趋势与实战策略】}

		      \textbf{考场秒杀}：类=模板，对象=实例
	      \end{bluebox}

	      \vspace{1em}

	\item \textbf{以下哪个是Java中的对象引用传递？}
	      \begin{enumerate}[label=\Alph*.]
		      \item 基本数据类型参数传递
		      \item 对象作为参数传递
		      \item 数组作为参数传递
		      \item 以上都是
	      \end{enumerate}

	      \begin{bluebox}{答案与深度解析}
		      \textbf{答案：B. 对象作为参数传递}

		      \textbf{【核心认知框架>}
		      \begin{itemize}[label={}, leftmargin=*, nosep]
			      \item \textbf{题目本质}：考查Java中参数传递机制——值传递 vs 引用传递
			      \item \textbf{关键区分点}：Java只有值传递，对象传递的是引用的副本（值）
		      \end{itemize}

		      \textbf{【选项原理深度拆解】}
		      \begin{itemize}[label={}, nosep]

		      \item \textbf{基本数据类型参数传递 — 错误（是值传递）}
		      \begin{itemize}[label={}, nosep]
			      \item 传递的是值的副本
			      \item 形参的改变不影响实参
		      \end{itemize}

		      \item \textbf{对象作为参数传递 — 正确（引用副本传递）}
		      \begin{itemize}[label={}, nosep]
			      \item 所有参数传递都是值传递
			      \item 对象传递的是引用地址的副本
			      \item 通过形参可以访问堆中对象
			      \item 但不能改变实参引用指向
		      \end{itemize}

		      \item \textbf{数组作为参数传递 — 部分正确（但非最佳答案）}
		      \begin{itemize}[label={}, nosep]
			      \item 数组也是对象（引用类型）
			      \item 传递机制与对象相同
		      \end{itemize}

		      \item \textbf{以上都是 — 错误（基本类型不是）}
		      \end{itemize}

		      \textbf{【应背诵知识点>}
		      \begin{itemize}[label={}, nosep]
			      \item 基本类型：值传递（传副本）
			      \item 引用类型：传引用地址的副本
			      \item 注意：不能改变实参引用指向
			      \item Java没有真正的引用传递
		      \end{itemize}

		      \textbf{【命题趋势与实战策略】}

		      \textbf{考场秒杀}：Java只有值传递，对象传引用副本
	      \end{bluebox}

\end{enumerate}

\end{document}
